\section{A mixed cubic, a mixed quartic, and two short convergent identities}
\label{mtp:examples}

In this section a subscript lists the tail orders: $H_{224}$ means
$H_{(2,2,4)}$, whereas $H_4$ is the generator of one fourth-order tail.
Put
\begin{equation}\label{mtp:E5}
 E_5=\frac{15}{2}\zeta(5)-3\zeta(2)\zeta(3)
     =\zeta(2,3)+\zeta(4,1)+\zeta(5).
\end{equation}
The second expression follows from the classical Euler reductions for
odd-weight double zeta values \cite[Equations (4), (9)]{Hoffman}.
The lower generators supplied by Theorem~\ref{mtp:generator} are
\begin{align}
 H_2&=-\frac{t^2}{4},
 &H_4&=\frac{\zeta(2)t^2}{2}-\frac{t^4}{48},\label{mtp:single-generators}\\
 H_{22}&=3\zeta(3)t+2I^2\bigl(\cosh(t/2)G_2\bigr),\notag\\
 H_{24}&=E_5t+\frac{\zeta(3)t^3}{4}
       +I^2\bigl(\cosh(t/2)G_4\bigr)
       +I^4\bigl(\cosh(t/2)G_2\bigr),\label{mtp:lower-generators}\\
 H_{222}&=3I^2\bigl(\cosh(t/2)G_{22}\bigr)
      -\frac{17}{8}\zeta(4)t^2-\frac{\zeta(2)t^4}{48}
      +\frac{t^6}{2880}.\notag
\end{align}
As throughout, $G_J=H_J/(2\sinh(t/2))$. These lower formulas are
inherited special cases, included to specify the subsequent mixed
formulas without unevaluated initial data.

\subsection{All moments of \texorpdfstring{$T_2^2T_4$}{T2 squared times T4}}

\begin{theorem}[A six-coordinate mixed cubic]
\label{mtp:triple}
The complete mixed-cubic generator is
\begin{equation}\label{mtp:H224}
 \boxed{\begin{aligned}
 H_{224}(t)={}&2I^2\bigl(\cosh(t/2)G_{24}\bigr)
                 +I^4\bigl(\cosh(t/2)G_{22}\bigr)\\
       &-\frac5{16}\zeta(6)t^2-\frac7{96}\zeta(4)t^4
                    -\frac{\zeta(2)t^6}{1440}+\frac{t^8}{161280}.
 \end{aligned}}
\end{equation}
Consequently every odd moment $M_{224,2r+1}$ belongs to the rational
span of
\begin{equation}\label{mtp:triple-coordinates}
                         1,\ \zeta(2),\ \zeta(3),\ \zeta(4),\ E_5,\ \zeta(6).
\end{equation}
The first six odd moments have full rank on these six formal labels.
Thus every later odd moment is a rational linear combination of
$M_{224,1},M_{224,3},\ldots,M_{224,11}$.
\end{theorem}

\begin{proof}
Apply \eqref{mtp:generator-formula} to the three labeled orders
$(2,2,4)$. Its nonpolynomial terms are
\[
 2I^2(\cosh G_{24})+I^4(\cosh G_{22})
                      -\frac16I^4H_4-\frac13I^6H_2.
\]
Here and only in this intermediate line, $\cosh$ abbreviates
$\cosh(t/2)$. The anchor polynomial from \eqref{mtp:R} is
\[
          -\frac5{16}\zeta(6)t^2-\frac7{96}\zeta(4)t^4
                                      -\frac{\zeta(2)t^6}{2160}.
\]
For example, its quadratic coefficient simplifies using
\[
 \mathcal S_{\{2,2\},2}
  =4\zeta(2,2,2)+2\zeta(4,2)+2\zeta(2,4)+\zeta(6),
\]
$\zeta(2)\zeta(4)=\zeta(2,4)+\zeta(4,2)+\zeta(6)$, and
$\zeta(2,2,2)=3\zeta(6)/16$. The quartic coefficient uses
$\zeta(2)^2=5\zeta(4)/2$. Inserting the two singleton generators
\eqref{mtp:single-generators} gives exactly \eqref{mtp:H224}.
The lower formulas \eqref{mtp:lower-generators} then prove the asserted
six-coordinate span for every Taylor coefficient.

In the displayed coordinate order, the first six odd moments have
coefficient matrix
\begin{equation}\label{mtp:triple-matrix}
\begingroup\setlength{\arraycolsep}{3pt}\renewcommand{\arraystretch}{1.3}
\begin{pmatrix}
0&0&0&0&1&-5/16\\
0&1/12&1&-7/16&-1/6&5/64\\
-19/360&-29/360&-7/12&35/96&47/720&-35/768\\
761/20160&2257/20160&61/120&-343/768&-257/5040&155/3072\\
-967/25200&-71059/302400&-2419/3360&217/256&9521/134400&-381/4096\\
298871/3193344&3264013/4561920&128981/80640&-9779/4096&-40063/253440&12775/49152
\end{pmatrix}.
\endgroup
\end{equation}
Its exact rational determinant is
\begin{equation}\label{mtp:triple-determinant}
                         -\frac{309971}{13243436236800}\ne0.
\end{equation}
The entries follow by finite Taylor extraction from
\eqref{mtp:H224}; ordinary rational elimination proves
\eqref{mtp:triple-determinant}. Both operations are independently
replayed in the supplied verifier. Inverting this finite matrix proves
the reconstruction assertion. It asserts rank on specified formal
coordinates, not arithmetic independence of the evaluated constants.
\end{proof}

For the first and the first subtracted moments, respectively, this gives
\begin{align}
 \sum_{n\ge0}x_nT_2(n)^2T_4(n)
  &=E_5-\frac5{16}\zeta(6),\label{mtp:triple-m1}\\
 \sum_{n\ge0}\left\{x_n^5T_2(n)^2T_4(n)-\frac13\right\}
  &=-\frac{19}{360}-\frac{29}{360}\zeta(2)-\frac7{12}\zeta(3)
       +\frac{35}{96}\zeta(4)\notag\\
  &\qquad+\frac{47}{720}E_5-\frac{35}{768}\zeta(6).
       \label{mtp:triple-m5}
\end{align}
The subtraction in the second line is fixed by
\[
 T_2(n)^2T_4(n)=\frac1{3x_n^5}-\frac2{9x_n^7}
                    +\frac{211}{1080x_n^9}-\frac{337}{1260x_n^{11}}
                    +O(x_n^{-13}).
\]
In polygamma notation its product is
$\psi'(n+1)^2\psi'''(n+1)/6$.

\begin{corollary}[A short cubic identity]
\label{mtp:short-cubic}
The following series is ordinary and absolutely convergent:
\begin{equation}\label{mtp:short-cubic-formula}
 \boxed{\sum_{n=0}^{\infty}T_2(n)T_4(n)
                     \bigl(1-x_nT_2(n)\bigr)=\frac5{16}\zeta(6).}
\end{equation}
Its summand is $O(x_n^{-6})$.
\end{corollary}

\begin{proof}
The classical unweighted two-tail evaluation gives
$\sum_{n\ge0}T_2(n)T_4(n)=E_5$; it follows directly by sorting the
two indices and is also the linear coefficient in
\eqref{mtp:lower-generators}. Subtract \eqref{mtp:triple-m1}.
Since $T_2(n)=x_n^{-1}-(12x_n^3)^{-1}+O(x_n^{-5})$, the stated
decay follows directly.

There is also a finite telescoping proof, independent of the moment
recurrence. Define
\[
 J(N)=-\frac{N^2}{2}T_2(N)^2T_4(N)+\frac16T_2(N)^3
            -\frac12T_2(N)T_4(N)-\frac16T_6(N).
\]
Direct expansion using $T_p(N-1)=T_p(N)+N^{-p}$ gives the exact identity
\begin{equation}\label{mtp:cubic-telescoper}
 J(N)-J(N-1)=T_2(N-1)T_4(N-1)
                    \bigl(1-(N-1/2)T_2(N-1)\bigr),\qquad N\ge1.
\end{equation}
Every term in $J(N)$ tends to zero. At the other endpoint,
\[
 J(0)=\frac16\zeta(2)^3-\frac12\zeta(2)\zeta(4)-\frac16\zeta(6)
             =-\frac5{16}\zeta(6).
\]
Thus finite summation of \eqref{mtp:cubic-telescoper} gives
\eqref{mtp:short-cubic-formula}. In particular, neither a conjectural
arithmetic relation nor a regularized limiting argument enters this
short evaluation.
\end{proof}

A purely rational identity is also available. Put
\begin{align*}
 P(x)=x\bigl(&2526903+80590275x^2+474388272x^4\\
             &+753887520x^6+353617920x^8+42197760x^{10}\bigr),\\
 Q(x)={}&14065920x^6+108495360x^4+180958272x^2+\frac{338793712}{7}.
\end{align*}
Then
\begin{equation}\label{mtp:rational-sum}
 \sum_{n=0}^{\infty}\bigl\{P(x_n)T_2(n)^2T_4(n)-Q(x_n)\bigr\}
                                      =-6199420.
\end{equation}
The six integer coefficients of $P$ multiply the matrix
\eqref{mtp:triple-matrix} to $(-6199420,0,0,0,0,0)$.
The displayed centered asymptotic coefficients give exactly $Q$;
the residual summand is $O(x_n^{-2})$. Thus this identity follows from
finite rational arithmetic and an ordinary convergent sum.

\subsection{All moments of \texorpdfstring{$T_2^3T_4$}{T2 cubed times T4}}

Define the convergent, weight-nine anchor
\begin{equation}\label{mtp:A9}
 A_9=\sum_{n=0}^{\infty}T_2(n)^3T_4(n)
   =\sum_{\pi\models\{1,2,3,4\}}
        \zeta(P_{B_1},\ldots,P_{B_{\ell-1}},P_{B_\ell}-1),
 \qquad(p_1,p_2,p_3,p_4)=(2,2,2,4).
\end{equation}
The ordered-partition sum is finite; its grouped fifteen-term expansion
is included in the exact verification record. The equality is the
classical minimum-index tail identity, with the $n=0$ term included.
Set
\begin{equation}\label{mtp:B7}
 B_7=\frac{59}{24}\zeta(3)\zeta(4)
              +\frac{23}{12}\zeta(2)\zeta(5)-\frac{203}{32}\zeta(7).
\end{equation}

\begin{theorem}[An eight-coordinate mixed quartic]
\label{mtp:quartic}
The complete even-moment generator is
\begin{equation}\label{mtp:H2224}
 \boxed{\begin{aligned}
 H_{2224}(t)={}&A_9t+B_7t^3-\frac{E_5t^5}{180}\\
 &+3I^2\bigl(\cosh(t/2)G_{224}\bigr)
       +I^4\bigl(\cosh(t/2)G_{222}\bigr)\\
 &-\frac12I^4H_{24}-\frac12I^6H_{22}.
 \end{aligned}}
\end{equation}
Every even moment is therefore in the rational span of the eight
specified coordinates
\begin{equation}\label{mtp:quartic-coordinates}
                1,\ \zeta(2),\ \zeta(3),\ \zeta(4),\ \zeta(6),\ E_5,\ B_7,\ A_9.
\end{equation}
The first eight even moments, of orders $0,2,\ldots,14$, have full
formal rank eight. Their coefficient determinant in this order is
\begin{equation}\label{mtp:quartic-determinant}
                      -\frac{806331520717}{2249069145523200000}\ne0.
\end{equation}
Every later even moment is an explicit rational linear combination
of these eight moments.
\end{theorem}

\begin{proof}
There are three singleton deletions of order two and one of order four.
Among the even proper subsets there are three pairs of type $(2,2)$
and three of type $(2,4)$; multiplication by $B_2=1/6$ gives the two
coefficients $1/2$ in \eqref{mtp:H2224}. The anchor polynomial has
degrees one, three, and five. Its linear coefficient is $A_9$.
Its cubic coefficient is the explicit finite combination
\begin{align}
 B_7={}&\zeta(2,2,2,1)+\frac12\zeta(4,2,1)+\frac12\zeta(2,4,1)
          +\frac12\zeta(2,2,3)\notag\\
       &+\frac16\zeta(6,1)+\frac16\zeta(4,3)-\frac1{12}\zeta(7).
       \label{mtp:B7-MZV}
\end{align}
For completeness, let
\[
 Q_7=\sum_{n\ge0}T_2(n)^4
     =51\zeta(3)\zeta(4)+10\zeta(2)\zeta(5)-\frac{301}{4}\zeta(7).
\]
This is Hoffman's classical quartic evaluation, after adding the
$n=0$ term to his convention \cite[following Corollary 4.4]{Hoffman}.
Sorting its four tail indices shows
\[
 24B_7=Q_7-4\zeta(2,5)-2\zeta(4,3)-3\zeta(7).
\]
Euler's double-zeta evaluations
\[
 \zeta(2,5)=10\zeta(7)-4\zeta(2)\zeta(5)-2\zeta(3)\zeta(4),
 \quad \zeta(4,3)=17\zeta(7)-10\zeta(2)\zeta(5)
\]
give \eqref{mtp:B7}. Finally, the quintic anchor coefficient is
$-(2\zeta(2,3)+\zeta(5))/240=-E_5/180$.
This proves the generator identity and the eight-coordinate span.

The first eight coefficient rows are obtained by Taylor extraction
from the displayed finite system
\eqref{mtp:single-generators}, \eqref{mtp:lower-generators},
\eqref{mtp:H224}, and \eqref{mtp:H2224}. The exact record gives all
$64$ rational entries; the standalone verifier constructs them and
computes \eqref{mtp:quartic-determinant} by rational elimination.
This is a finite nonvanishing certificate. Since the generator already
proves the span for every order, this one exact determinant proves
reconstruction of every later moment from the stated prefix.
No independence of the actual period values is required.
\end{proof}

The first three even moments, in compact form, are
\begin{align}
 M_{2224,0}&=A_9,\notag\\
 M_{2224,2}&=E_5-\frac5{16}\zeta(6)+2B_7-\frac{A_9}{12},
       \label{mtp:quartic-first}\\
 M_{2224,4}&=\frac{\zeta(2)}{20}+\frac32\zeta(3)-\frac{11}{16}\zeta(4)
             +\frac1{16}\zeta(6)-\frac{23}{60}E_5-B_7+\frac7{240}A_9.
       \notag
\end{align}
All three are ordinary convergent sums without subtraction, since
$T_2^3T_4=O(x_n^{-6})$.

\begin{corollary}[A short mixed quartic--cubic identity]
\label{mtp:short-quartic}
One has the ordinary convergent evaluation
\begin{equation}\label{mtp:short-quartic-formula}
 \boxed{\begin{aligned}
 &\sum_{n=0}^{\infty}T_2(n)^2T_4(n)
       \left\{\left(x_n^2+\frac1{12}\right)T_2(n)-x_n\right\}\\
 &\qquad=\frac{59}{12}\zeta(3)\zeta(4)
                +\frac{23}{6}\zeta(2)\zeta(5)-\frac{203}{16}\zeta(7).
 \end{aligned}}
\end{equation}
The summand is $O(x_n^{-8})$.
\end{corollary}

\begin{proof}
The series on the left is
$M_{2224,2}+M_{2224,0}/12-M_{224,1}$. Equations
\eqref{mtp:triple-m1} and \eqref{mtp:quartic-first} reduce it to $2B_7$,
which is the displayed right side. The centered expansion gives
\[
 \left(x^2+\frac1{12}\right)T_2(n)-x
       =\frac1{45x^3}+O(x^{-5}),\qquad x=x_n.
\]
Multiplication by $T_2^2T_4=O(x^{-5})$ proves the asserted decay.
\end{proof}

\subsection{Independent verification and further questions}

The exact verifier uses two different constructions. The first
implements Theorem~\ref{mtp:recurrence}. The second starts from
ordinary finite summation by parts. If
\[
 \mathcal T_{I,q}^{\mathrm{cut}}
  =\operatorname{CT}_N\sum_{n=1}^Nn^qF_I(n),\qquad
       \sum_{n=1}^Nn^q=\sum_a f_{q,a}N^a,
\]
then for $q\ge0$ its independent recurrence is
\begin{equation}\label{mtp:independent-cutoff}
 \mathcal T_{I,q}^{\mathrm{cut}}
   =\sum_a f_{q,a}b_I(a)
    +\sum_{\varnothing\ne S\subseteq I}
       \left\{\sum_a f_{q,a}
                    \mathcal T_{I\setminus S,a-P_S}^{\mathrm{cut}}
                    -\mathcal T_{I\setminus S,q-P_S}^{\mathrm{cut}}\right\}.
\end{equation}
For $q<0$ and nonempty $I$, this is the convergent anchor
$\mathcal T_{I,-q}$. The empty-product cases are zero for $q\ge0$,
Euler's constant for $q=-1$, and $\zeta(-q)$ for $q\le-2$.
The finite identity uses $F_q(n-1)=F_q(n)-n^q$, so its cutoff constants
have exactly the required convention.

For all multisets of one through four orders in $\{2,4,6\}$, the two
routes agree in $306$ regular moments through order $17$, coefficient
by coefficient on unreduced MZV words. Thus no unproved MZV relation
is used to make the comparison succeed. Every Euler-constant label
cancels. There are another $825$ exact boundary-parity checks,
$29$ checks of the simplified mixed generators, the six- and
eight-coordinate determinant certificates, and the rational row in
\eqref{mtp:rational-sum}. The finite cubic telescoper is checked
independently as an exact identity in four Laurent-polynomial variables.
The verifier needs only the Python standard
library; its complete rational outputs accompany the article.

The separate numerical diagnostic evaluates the ordinary sums at
$90$ decimal digits, with $40$ or $64$ initial terms and fixed centered
asymptotic tail completions. For the four displayed cubic moments the
$40$-term residuals are between $3.4\times10^{-72}$ and
$1.6\times10^{-62}$; the $64$-term residuals are below $10^{-76}$.
The quartic second-moment relation has residual below $3.4\times10^{-72}$,
and the larger rational identity has residual below $1.3\times10^{-65}$.
These are numerical diagnostics, not interval certificates. The proofs
are the finite identities, the analytic continuation argument, and
exact rational elimination.

Several precise continuations remain:
\begin{enumerate}
\item Determine the smallest useful formal arithmetic space for a
      general multiset of even orders, and find a bound for the moment
      prefix needed to recover it. The examples above have dimensions
      six and eight, far below the unrestricted weight-based list.
\item Find further low-degree combinations of mixed moments whose
      largest-weight anchors cancel, as in
      \eqref{mtp:short-cubic-formula} and \eqref{mtp:short-quartic-formula}.
      A systematic classification could produce ordinary convergent
      identities with only single zeta values and their products.
\item Extend the boundary polynomial to common noninteger shifts.
      The shifted lattice no longer annihilates every nonnegative even
      power, so the finite endpoint correction must be retained.
\item Treat odd tail orders or spectral-order derivatives. Their
      centered parity and logarithmic asymptotic terms differ from the
      hypotheses proved here; simple polarization of the present
      theorem does not supply those cases.
\item Identify minimal functional multiple-polylogarithm depth under
      known functional relations. The bound $d-1$ is constructive,
      and makes no assertion that a smaller expression is impossible.
\end{enumerate}
